\chapter{Proposal Feedback and Iteration Plan}
\label{ch:plan}
\section{Proposal Feedback Summary}
The proposal feedback and its resolution are recorded in the action register below and in Appendix~\ref{app:actions}. \todo{insert the actual proposal-defense feedback received from the evaluators and supervisor}

The scope decisions already reflected in the project (descoping multi-camera re-identification, shelf-level interaction detection, and long-term layout evaluation) came out of the feasibility triage described in Chapter~\ref{ch:research} and are tracked here as actions.

\section{Proposal Action Register}
\begin{longtable}{P{0.06\textwidth}P{0.28\textwidth}P{0.25\textwidth}P{0.14\textwidth}P{0.17\textwidth}}
\toprule
No. & Feedback / Required Change & Action Taken & Status & Evidence\\
\midrule
\endhead
1 & \todo{evaluator feedback} & Scope narrowed to a single-store, three-camera, CPU-only prototype after a feasibility triage. & Done & Chapter~\ref{ch:intro} scope\\
2 & \todo{evaluator feedback} & Full multi-camera identity linking descoped. An experimental two-camera handoff MVP built instead. & Done & Chapter~\ref{ch:iter1}\\
3 & \todo{evaluator feedback} & POS analytics added with deterministic metrics and cited assistant answers. & Done & Chapter~\ref{ch:iter1}\\
4 & \todo{evaluator feedback} & Privacy approach fixed as non-biometric, and results framed as decision support. & Done & Chapter~\ref{ch:intro}\\
\bottomrule
\end{longtable}

\section{Iteration Strategy}
Development is iterative, with each iteration delivering a working increment that is demonstrated and evaluated before the next begins. Iteration 1 builds the platform foundation: the web application, the POS analytics service, secure access, and the camera-handoff MVP. Iteration 2 connects the vision agents to the platform through the orchestrator and evaluates them on store footage. Work is tracked on feature branches in the project repository. \todo{confirm iteration boundaries and dates with the supervisor}

\section{Iteration 1 Plan}
\textbf{Objectives:} establish the application shell and data flows that later analytics plug into, and prove the riskiest vision idea (cross-camera handoff) early.

\textbf{Selected features:} manager and employee pages; role-aware authentication; reports with PDF export; alerts feed; system status page; POS import, metrics, forecasts, recommendations, and cited assistant; camera-handoff pipeline and its web endpoint.

\textbf{Deliverables:} running frontend, local backend (\texttt{/pos/*} and \texttt{/camera/*}), handoff JSON output, and demonstration evidence (Chapter~\ref{ch:iter1}).

\section{Iteration 2 Plan}
\textbf{Objectives:} replace prototype data with real agent output and evaluate it.

\textbf{Selected features:} Detection, Classification, Engagement (including interaction factors), and Trajectory agents; the LangGraph orchestrator; live data in the dashboard pages; the Insight Narration agent; and quantitative evaluation on labelled clips.

\textbf{Deliverables:} integrated pipeline, evaluation report, and updated dashboard bound to real outputs (Chapter~\ref{ch:fyp2plan}).

\section{Work Breakdown and Responsibilities}
\begin{table}[H]
\centering
\caption{Planned work breakdown}
\begin{tabular}{P{0.3\textwidth}P{0.45\textwidth}P{0.15\textwidth}}
\toprule
Work package & Description & Owner\\
\midrule
Frontend and UX & Next.js pages, design system, charts, PDF export & \todo{name}\\
POS analytics backend & FastAPI, SQLite, forecasts, recommendations, assistant & \todo{name}\\
Vision pipeline & Detection, tracking, classification, engagement, trajectory & \todo{name}\\
Camera handoff & Boundary events, time window, appearance matching & \todo{name}\\
Orchestration and narration & LangGraph graph, reconciliation, narration & \todo{name}\\
Testing and documentation & Test cases, evaluation, report & \todo{name}\\
\bottomrule
\end{tabular}
\end{table}
The team consists of Moiz Ansari, Saqib Atif, and Ruhab Ahmad. Individual ownership is recorded in Chapter~\ref{ch:mgmt}.

\section{Progress Against Plan}
\begin{table}[H]
\centering
\caption{Planned versus actual progress at mid-report}
\begin{tabular}{P{0.34\textwidth}P{0.2\textwidth}P{0.3\textwidth}}
\toprule
Milestone & Planned & Actual\\
\midrule
Web application (manager and employee) & Iteration 1 & Implemented\\
POS analytics service & Iteration 1 & Implemented (local)\\
Camera-handoff MVP & Iteration 1 & Implemented (experimental)\\
Floor-distance CLI & Current prototype increment & Implemented provisionally\\
Vision agents wired to dashboard & Iteration 2 & In progress (prototype data)\\
LangGraph orchestration & Iteration 2 & In progress\\
Independent physical / labelled accuracy evaluation & Iteration 2 & Pending; floor software checks already executed\\
\bottomrule
\end{tabular}
\end{table}

\section{Delays, Deviations and Recovery Plan}
\textbf{Deviation 1.} Full cross-camera re-identification was descoped to a conservative handoff MVP because of the CPU-only constraint and limited footage. \textbf{Deviation 2.} A POS analytics pillar was implemented as a local FastAPI and SQLite service ahead of the vision integration, so the dashboard has real data in one area while the vision agents mature. \textbf{Residual gap.} The dashboard's camera-derived metrics still use prototype data. \textbf{Recovery:} the first FYP-2 sprint binds the Detection and Engagement outputs to the dashboard, and quantitative evaluation is scheduled for the middle of the semester so that tuning time remains.
